\chapter{个人体会与收获 在巨变的时代中定位}

\section{知识体系的整合与重构}

通过本学期前沿讲座的学习，我最大的收获之一是原有的相对零散的计算机知识体系得到整合与重构。此前我学习操作系统、计算机网络、计算机组成原理、人工智能等课程时各门课相对独立。而本次讲座中，张珊教授从网络基本问题出发引出对存算传融合的需求，李波教授从高性能计算角度展示处理器、内存、网络互连如何协同，黄迪教授展示视觉算法如何在移动端和边缘端受限于算力和带宽，刘偲教授展示具身智能如何依赖高效的通信。这些内容将我脑海中的知识孤岛连接成一片完整的图景。我开始理解为什么操作系统中的虚拟内存管理要与网络协议栈中的RDMA结合，为什么计算机体系结构中的多核缓存一致性会影响分布式机器学习的效率，为什么计算机视觉中的模型压缩直接影响自主无人系统的响应速度。这种知识融合的视角将是我未来从事系统型研究的重要基础。

\section{科研思维与方法论的提升}

李波老师的科研方法论讲座对我影响最为深远。此前我常以为科研就是读论文加做实验加写论文的线性过程。李老师指出研究的核心在于搜索与比较，搜索所有可能的方法并比较，比较的结论引导新的发现。这让我意识到文献综述不仅仅是罗列前人工作，而是要通过深度比较发现现有方法的不足，从而找到自己的创新点。他还强调计算机科学研究的最难点是提出自己的创新点，而创新往往受限于基础科学的突破。这让我明白作为本科生虽然很难做出基础科学层面的突破，但可以在已有理论的基础上通过迁移、组合、改进现有方法在应用层面做出有价值的创新。关于论文写作，李老师强调数学模型往往是本质性的，这提醒我在今后的研究中要重视数学推导，而不仅仅满足于实验结果的漂亮。

\section{对人工智能未来发展的判断}

通过李舟军老师和刘偲老师的讲座，我对人工智能的未来发展有了更清醒的判断。通用大模型如DeepSeek-V3和GPT-4将作为基础设施存在，类似今天的操作系统或数据库，而真正的商业价值将爆发在通专结合的垂直领域，即那些需要深厚领域知识、高可靠性、实时性的场景如医疗诊断、工业控制、金融风控。这也意味着作为计算机专业的学生，如果仅仅掌握调用大模型应用程序接口的技能，未来很容易被替代。相反，如果能够深入某个垂直领域如生物信息、智能制造、航空航天，成为既懂人工智能又懂领域的复合型人才，将具备持久的竞争力。具身智能将是人工智能下一个十年的主战场，但其技术难度远超纯软件人工智能，需要硬件、算法、系统、通信等多方面的协同突破。我对刘偲老师团队在无人机视觉语言动作模型方面的工作非常感兴趣，希望未来有机会在该方向深入探索。

\section{对个人学习路径的规划}

结合以上收获，我对未来两年的学习与研究方向有了初步规划。

\textbf{夯实基础。} 继续深入学习操作系统、计算机组成原理、计算机网络、算法设计与分析等核心课程，尤其要理解底层原理，避免浮于表面。

\textbf{拓展交叉领域。} 选修或自学一些与人工智能加系统、人工智能加机器人相关的课程，如机器人学、控制理论、边缘计算等。同时关注生物信息、智慧城市等应用领域，培养跨学科思维。

\textbf{积极参与科研实践。} 争取进入实验室跟随老师或研究生参与实际项目。刘偲老师、黄迪老师、张珊老师等课题组的研究方向都非常前沿，我希望能够加入其中，从具体问题入手，锻炼发现问题和解决问题的能力。

\textbf{培养批判性思维。} 在阅读论文和聆听报告时有意识地追问这个方法的核心贡献是什么，它的局限在哪里，能否改进或迁移到其他场景。避免盲目崇拜顶会论文或工业界巨头的工作。

\textbf{关注国产技术生态。} 李舟军老师提到DeepSeek与华为昇腾的合作，张珊老师提到自主计算对国产芯片的意义。我认识到在当前国际环境下掌握自主可控的关键技术至关重要。我会关注国产操作系统如OpenHarmony、国产人工智能芯片如昇腾和寒武纪以及国产数据库和中间件，为将来投身国家信息技术创新事业打下基础。

\section{课程带来的职业与人生启示}

这次课程也让我对未来的职业选择有了更深的思考。李波老师提到核心技术和非核心技术的差异在于不可替代性。虽然短期内从事应用开发可能更容易找到工作，但长期来看掌握核心技术的人才才能真正立于不败之地。周忠教授展示的虚拟现实与增强现实在医疗、航天等领域的应用，刘偲教授展示的无人机在应急救援、物流配送中的价值，让我感受到技术从业者的社会责任感。我们不仅仅是在写代码和调参数，更是在用技术解决实际问题、改善人们的生活、甚至保卫国家安全。这种使命感将激励我在未来的学习和工作中始终牢记将论文写在祖国大地上。